% !TITLE 登高
% !AUTHOR 杜甫
% !DYNASTY 唐
% !ORDER 30

\begin{poem}
风急天高猿啸哀，渚清沙白鸟飞回\textsuperscript{1}。
无边落木萧萧下\textsuperscript{2}，不尽长江滚滚来。
万里悲秋常作客\textsuperscript{3}，百年多病独登台\textsuperscript{4}。
艰难苦恨繁霜鬓\textsuperscript{5}，潦倒新停浊酒杯\textsuperscript{6}。
\end{poem}

\begin{annotations}
\notetext{1}{渚清沙白鸟飞回：水中小洲清澈，沙滩洁白，鸟儿盘旋飞回。渚（zhǔ），水中小块陆地。六个意象并排——风、天、猿、渚、沙、鸟，像一幅工笔山水。}
\notetext{2}{落木：落叶。萧萧，落叶的声音。无边无际的落叶纷纷飘落。}
\notetext{3}{万里悲秋常作客：漂泊万里，在秋天里悲伤，常年做客他乡。“万里”写空间之远，“常”写时间之久。}
\notetext{4}{百年多病独登台：人生百年，疾病缠身，独自登上高台。“百年”写人生之短暂，“多病”写身体之衰残。}
\notetext{5}{繁霜鬓：白发像霜一样繁密。艰难和苦恨的累积，让白发愈加浓厚。}
\notetext{6}{潦倒新停浊酒杯：潦倒困顿，最近又戒了酒。新停，刚刚停止。连借酒消愁的最后一招也用不了了。}
\end{annotations}

\begin{appreciation}
《登高》被明代学者胡应麟推为"古今七律第一"。七律是唐代律诗里技术要求最高的体裁：要平仄，要对仗，要起承转合。这首诗拿的是这个项目的历史最高分，至今没人打破。

风急天高猿啸哀，渚清沙白鸟飞回——十四个字，六个意象：风、天、猿、渚、沙、鸟。风是急的，天是高的，猿的叫声是悲哀的；小洲是清的，沙滩是白的，鸟在盘旋。没有一个废字。写这首诗时杜甫在夔州，长江边上。三峡的猿啼出了名的哀——当年李白写"两岸猿声啼不住"是归心似箭，杜甫听来句句都是悲音。

无边落木萧萧下，不尽长江滚滚来——落木就是落叶，但"木"比"叶"多一份重量。萧萧是落叶的声音，滚滚是江水奔流的样子，一个是听觉，一个是视觉，都配着"无边""不尽"的浩大。这就是对仗：两句像照镜子，字字严丝合缝。更要紧的是意境：落叶无边，是时光在流逝；江水不尽，是人生太短暂。一个人站在高处，看见无穷的落和无穷的流，会觉得自己特别渺小。

万里悲秋常作客，百年多病独登台——宋人罗大经数过，这十四个字里有八层悲：万里离家，悲秋，作客，常年作客，百年苦短，多病，登台添愁，最后独登台——别人重阳登高是一家人，他只有一个人。八层悲一层压一层，最后全压在"独"字上。

艰难苦恨繁霜鬓，潦倒新停浊酒杯——繁霜鬓，白发像霜一样密。苦恨累积到一定程度，会在头发上显出来。新停，是最近刚刚戒了酒。为什么戒？病了。酒是杜甫唯一的消遣，如今连这最后一杯也喝不成了——一个人的慰藉被剥夺到最后，什么都不剩。

重阳本是登高节日，插茱萸、喝菊花酒、全家出门，而杜甫在夔州，孤身登上一座高台。中国的诗人写秋悲，从宋玉的"悲哉秋之为气也"到杜甫这一首，算是写到了头。而这悲的根，还是《采薇》里"我心伤悲，莫知我哀"，还是《离骚》里屈原的忧国忧民——个人的悲哀连着天下的动荡，这是中国诗的传统，杜甫把它推到了顶点。

写这首诗的第三年，杜甫死在湘江的船上。他写的时候大概不会想到，后世会把它推为千古第一律诗。他只是一个人，过了一个节，把眼前所见、心里所想，如实写了下来。

妹妹，这首诗哥哥建议你背下来。不是为了考试，是为了将来有一天你一个人在外地过节，会想起一千多年前也有一个人，一个人过节，把日子过成了诗。
\end{appreciation}
